\begin{remark}[Topological order and length dependence]%
	\label{rem:mpdo_bnt_label_length_independent_topological}
	Starting from the same-length product law and the condition that the
	vector $m$ is an idempotent for the multiplication induced by $c^{(1)}$,
	in the particular case of natural-number structure coefficients
	independent of the length, one recovers all known non-chiral
	topologically ordered phases together with a description of their
	excitations \cite{Bultinck2017Anyons}. The RFP density operators
	correspond to the boundary theories of the RFP representatives of those
	phases, which yields a large family of non-trivial examples. The source
	asks whether there exist RFPs whose structure coefficients depend on the
	length \cite[Section~4.5]{Cirac2016MPDO_arXiv}. The example below gives
	such coefficients for its displayed BNT presentation.
\end{remark}

\begin{theorem}[An RFP with length-dependent BNT structure coefficients]%
	\label{thm:mpdo_length_dependent_rfp_example}
	\lean{MPOTensor.BondTwoSingletonBaseModel.%
		normalizedSingletonCoeffs_coeff_one_ne_coeff_two}
	\lean{MPOTensor.BondTwoSingletonBaseModel.%
		normalizedSingletonCoeffs_not_lengthIndependent}
	\lean{MPOTensor.BondTwoSingletonBaseModel.%
		exists_isRFPViaTS_not_lengthIndependent_bntCoefficients}
	\leanok
	\uses{def:mpdo_bnt_label_length_independent, thm:mpdo_bond_two_singleton_normalized_algebra, thm:mpdo_bnt_algebra_tensor_to_rfp}
	There is a tensor $M$ in CF that generates MPDOs, is an RFP, and
	has a vertical BNT presentation with one element and structure
	coefficient $c^{(L)}=(\sqrt{1/2})^L$.
	In particular, $c^{(1)}\ne c^{(2)}$, so these structure coefficients
	depend on the length. This statement concerns the exhibited BNT
	presentation; it does not assert that length dependence is invariant
	under every possible presentation of $M$.
\end{theorem}
\begin{proof}\leanok
	\uses{thm:mpdo_bond_two_singleton_ambient_ghz, thm:mpdo_bond_two_singleton_cpsv_form, thm:mpdo_bond_two_singleton_normalized_algebra, thm:mpdo_bond_two_singleton_tensor_clause, thm:mpdo_bnt_algebra_tensor_to_rfp}
	The bond-dimension-two tensor with a one-element vertical BNT
	generates positive operators on every nonempty chain, is in CF, and
	satisfies the displayed BNT multiplication law and the corresponding
	tensor decomposition. The reverse implication of
	Theorem~\ref{thm:mpdo_bnt_algebra_tensor_to_rfp} supplies the two
	tpCPMs required for an RFP. Finally,
	$\sqrt{1/2}\ne1/2=(\sqrt{1/2})^2$, which proves the length dependence.
\end{proof}

\begin{theorem}[Length dependence under every uniform positive rescaling]%
	\label{thm:mpdo_bnt_label_rescaling_stable_length_dependence}
	\lean{MPOTensor.RescalingStableLengthDependentRFP.A}
	\lean{MPOTensor.RescalingStableLengthDependentRFP.R}
	\lean{MPOTensor.RescalingStableLengthDependentRFP.%
		physTraceTransfer_R_idempotent}
	\lean{MPOTensor.RescalingStableLengthDependentRFP.oneLabelChi}
	\lean{MPOTensor.RescalingStableLengthDependentRFP.oneLabelCoeffs}
	\lean{MPOTensor.RescalingStableLengthDependentRFP.oneLabelCoeffs_coeff}
	\lean{MPOTensor.RescalingStableLengthDependentRFP.%
		oneLabelCoeffs_coeff_one_ne_coeff_two}
	\lean{MPOTensor.RescalingStableLengthDependentRFP.%
		oneLabelCoeffs_not_lengthIndependent}
	\lean{MPOTensor.RescalingStableLengthDependentRFP.oneLabelChiScaled}
	\lean{MPOTensor.RescalingStableLengthDependentRFP.rescaledCoeffs}
	\lean{MPOTensor.RescalingStableLengthDependentRFP.rescaledCoeffs_coeff}
	\lean{MPOTensor.RescalingStableLengthDependentRFP.%
		oneLabelCoeffs_rescaling_stable_not_lengthIndependent}
	\leanok
	\uses{def:mpdo_bnt_label_coefficient_family, def:mpdo_bnt_label_length_independent, def:mpo_physical_trace_transfer}
	Theorem~\ref{thm:mpdo_length_dependent_rfp_example} gives length-dependent
	structure coefficients for a specified BNT presentation, as asked in
	\cite[lines~995--999]{Cirac2016MPDO_arXiv}.
	Its diagonal matrix $\chi$ has the single entry $\sqrt{1/2}$,
	and rescaling that entry to $1$ makes the coefficient constant.
	The diagonal matrix constructed here has two distinct positive entries,
	and its trace powers remain length dependent under every uniform
	positive rescaling. Define four matrices in $\MN{2}$ by
	\begin{align}
		A^0 & =\frac45E_{00}, &
		A^1 & =\frac45E_{11}, &
		A^2 & =\frac35E_{01}, &
		A^3 & =\frac35E_{10}.
		\notag
	\end{align}
	Identify the physical index set $\{0,1,2,3\}$ with
	$\{0,1\}\times\{0,1\}$ and let $R$ be the bond-dimension-four MPO
	tensor obtained by exchanging physical and virtual indices in the
	matrices $B^{ab}=A^a\otimes\overline{A^b}$ and applying the indicated
	normalization:
	\begin{align}
		R^{pq}_{ab} & =\frac{25}{32}
		\left(A^a\otimes\overline{A^b}\right)_{p,q}.
		\notag
	\end{align}
	Here $(p,q)$ are the physical indices and $(a,b)$ are the bond
	indices. Its physical-trace transfer is the rank-one projection
	\begin{align}
		\mathcal T_R & =\frac{25}{32}\ketbra{t}{t}, &
		t & =\left(\frac45,\frac45,0,0\right),
		\notag
	\end{align}
	where $t_a=\tr(A^a)$. It satisfies $\mathcal T_R^2=\mathcal T_R$,
	the zero-correlation-length condition of
	\cite[Definition~4.2, lines~736--741]{Cirac2016MPDO_arXiv}.
	Define the positive diagonal matrix
	$\chi=\diag(1,\frac{7}{25})$ and its trace powers by
	\begin{align}
		c^{(L)} & =\tr(\chi^L)=1+\left(\frac{7}{25}\right)^L.
		\notag
	\end{align}
	Hence $c^{(1)}\ne c^{(2)}$, so these coefficients are not length
	independent. For every real number $s>0$, the rescaled matrix
	$s\chi=\diag(s,\frac{7s}{25})$ has trace powers
	\begin{align}
		c_s^{(L)} & =\tr\!\left((s\chi)^L\right)
		=s^L\left(1+\left(\frac{7}{25}\right)^L\right),
		\notag
	\end{align}
	which are again not length independent. Unlike the single-entry
	matrix of Theorem~\ref{thm:mpdo_length_dependent_rfp_example}, no
	uniform positive rescaling makes both diagonal entries equal to one.

	This conclusion concerns the displayed trace powers and every uniform
	positive rescaling of $\chi$. The unscaled coefficients are identified
	below as the structure coefficients of an explicit one-element vertical
	BNT presentation of $R$. No such realization is asserted here for the
	arbitrary rescaled coefficients. Positivity, CF, and the RFP property
	of $R$ are proved separately. This example is motivated by the
	discussion in \cite[lines~995--1010]{Cirac2016MPDO_arXiv}; it is not
	stated there.
\end{theorem}
\begin{proof}\leanok
	\uses{def:mpo_physical_trace_transfer, def:mpdo_bnt_label_coefficient_family, def:mpdo_bnt_label_length_independent}
	The $(a,b)$ entry of $\mathcal T_R$ is
	$\frac{25}{32}\tr(A^a)\tr(A^b)$, and
	$\sum_a\tr(A^a)^2=2\cdot(\frac45)^2=\frac{32}{25}$,
	so $\mathcal T_R^2=\mathcal T_R$. The formulas for $c^{(L)}$ and
	$c_s^{(L)}$ follow by summing the powers of the diagonal entries of
	$\chi$ and $s\chi$.
	Write $\lambda=\frac{7}{25}$. If $c_s$ were length independent,
	comparing lengths one and two and lengths two and three would give
	\begin{align}
		s & =\frac{1+\lambda}{1+\lambda^2},
		\notag\\
		s & =\frac{1+\lambda^2}{1+\lambda^3}.
		\notag
	\end{align}
	Hence $(1+\lambda)(1+\lambda^3)=(1+\lambda^2)^2$, that is,
	$\lambda(1-\lambda)^2=0$. This fails at $\lambda=\frac{7}{25}$.
\end{proof}

\begin{definition}[The \texorpdfstring{$\Z_2$}{Z2}-twisted quantum dimer]%
	\label{def:mpdo_twisted_dimer_tensor}
	\lean{MPOTensor.TwistedDimer.physIdx}
	\lean{MPOTensor.TwistedDimer.Cmat}
	\lean{MPOTensor.TwistedDimer.tau}
	\lean{MPOTensor.TwistedDimer.coef}
	\lean{MPOTensor.TwistedDimer.T}
	\lean{MPOTensor.TwistedDimer.block}
	\lean{MPOTensor.TwistedDimer.T_apply_blocks}
	\leanok
	Each site carries two qubits $L$, $R$ and a flag qubit $f$, so the
	physical index set is $\{0,1\}^3$, read as $(l,r,f)$. The bond index set
	is also $\{0,1\}^3$, read as $(p,p',k)$, where $k$ labels one of two
	horizontal blocks. Let $X$ and $Z$ denote the Pauli matrices.
	With $x=\frac78$, $y=\frac18$, $P_\pm=\frac12(\Id\pm X)$,
	and the $2\times2$ matrices $C_0=xP_++yP_-$ and
	$C_1=xP_+-yP_-$, the bond matrix for the pair of physical indices
	$((l,r,f),(l',r',f'))$ is
	\begin{align}
		M^{(l,r,f),(l',r',f')}
		 & =\sum_{k\in\{0,1\}}\frac12\,C_k[l,l']\,(\tau_k)_{ff}
		\delta_{ff'}\,E_{(l,l',k),(r,r',k)}.
		\notag
	\end{align}
	Here $\tau_0=\Id$ and $\tau_1=Z$. Thus each bond matrix is a sum of
	two matrix units, one in each block. The tensor is the direct sum
	along $k$ of the two bond-dimension-four blocks obtained by keeping
	one summand. This example is motivated by the question of
	\cite[lines~995--1010]{Cirac2016MPDO_arXiv}; it is not stated there.
\end{definition}

\begin{theorem}[Closures and physical-trace transfer of the twisted dimer]%
	\label{thm:mpdo_twisted_dimer_closures}
	\lean{MPOTensor.TwistedDimer.physClose1_T}
	\lean{MPOTensor.TwistedDimer.T_mul_T}
	\lean{MPOTensor.TwistedDimer.physClose2_T}
	\lean{MPOTensor.TwistedDimer.physTraceTransfer_block_one}
	\lean{MPOTensor.TwistedDimer.physTraceTransfer_block_zero}
	\leanok
	\uses{def:mpdo_twisted_dimer_tensor, def:mpo_physical_trace_transfer}
	For a bond matrix $X$, the one-site physical closure $M_1(X)$ of the
	twisted dimer has entries
	\begin{align}
		[M_1(X)]_{(l,r,f),(l',r',f')}
		&=\sum_k\frac12\,C_k[l,l']\,(\tau_k)_{ff} \delta_{ff'}
		X_{(r,r',k),(l,l',k)}.
		\notag
	\end{align}
	The two-site closure $M_2(X)$ vanishes unless the consecutive physical
	indices satisfy $r_1=l_2$ and $r'_1=l'_2$. When these conditions hold,
	its entry is the sum over a common block label $k$ of the product of the
	two coefficients times $X_{(r_2,r'_2,k),(l_1,l'_1,k)}$.
	The physical-trace transfer of the block $k=0$ is
	$\sum_{l,r}C_0[l,l]\,E_{(l,l),(r,r)}$, and that of the block $k=1$
	vanishes because $\tr\tau_1=0$. Thus the latter block has nilpotent
	physical-trace transfer in the sense used by
	\cite[Definition~4.7, lines~815--822]{Cirac2016MPDO_arXiv}.
\end{theorem}
\begin{proof}\leanok
	\uses{def:mpdo_twisted_dimer_tensor}
	The trace of a matrix unit against $X$ picks one entry of $X$, and the
	product of two matrix units is a matrix unit or zero according to whether
	the inner indices agree. Since the block label is part of the bond index,
	only common labels survive. Summing the diagonal physical indices over
	the flag value gives $\sum_f(\tau_k)_{ff}$, which is $2$ for $k=0$
	and $0$ for $k=1$.
\end{proof}

\begin{theorem}[Positivity of the twisted quantum dimer]%
	\label{thm:mpdo_twisted_dimer_ismpdo}
	\lean{MPOTensor.TwistedDimer.mpo_T_entry_formula}
	\lean{MPOTensor.TwistedDimer.mpo_T_eq_smul_hadamard}
	\lean{MPOTensor.TwistedDimer.T_isMPDO}
	\leanok
	\uses{def:mpdo, def:mpdo_twisted_dimer_tensor}
	The closed operator $\rho^{(N)}(M)$ of the twisted quantum dimer is
	positive semidefinite for every positive chain length $N$, so $M$
	generates MPDOs. Writing $l_n$, $r_n$, $f_n$ for the three bits at
	site $n$ of a physical string $\sigma$, and using primes for those
	of $\tau$, the matrix element is
	\begin{align}
		\rho^{(N)}(M)_{\sigma\tau}
		 & =\frac1{2^N}
		\Id[\text{cyclic bond match}](\sigma,\tau)
		\delta_{f f'}\left(\prod_n C_0[l_n,l'_n]
		             +(-1)^{\lvert f\rvert}\prod_n C_1[l_n,l'_n]\right).
		\notag
	\end{align}
	The indicator is one exactly when both $\sigma$ and $\tau$ match
	their bond bits around the ring, $\delta_{ff'}$ is one exactly when the two
	flag strings agree, and $\lvert f\rvert$ is the number of sites carrying
	flag one. Thus $\rho^{(N)}(M)$ is $2^{-N}$ times the Hadamard product
	of the rank-one bond-matching indicator with the pullback, along the
	flag and left bits, of the sum of the $N$-fold Kronecker powers of
	the two local four-by-four factors
	$g_k[(f,l),(f',l')]=\delta_{ff'}\,C_k[l,l']\,(\tau_k)_{ff}$.
	The Schur product theorem then gives positivity.
\end{theorem}
\begin{proof}\leanok
	\uses{def:mpdo_twisted_dimer_tensor}
	Each bond matrix is a sum of two matrix units, one per block label,
	so an ordered product keeps a common block label and vanishes unless
	consecutive physical indices match on their bond bits. The surviving
	product is a single matrix unit whose row is determined by the left
	bits at the first site and whose column is determined by the right
	bits at the last site. Taking the trace imposes the wraparound match
	and leaves the sum over the block label of the product of the one-site
	coefficients. Each coefficient is half of $g_k$ evaluated at the flag
	and left bits, so the product along the ring is $2^{-N}$ times an entry
	of the $N$-fold Kronecker power of $g_k$.

	Neither the sum nor the difference of the two Kronecker powers is a
	Kronecker power, so both are proved positive semidefinite together by
	induction on the chain length. Separating the last tensor factor gives
	\begin{align}
		g_0^{\otimes(N+1)}\pm g_1^{\otimes(N+1)}
		 & =\frac12\Bigl[
			           (g_0^{\otimes N}+g_1^{\otimes N})\otimes(g_0\pm g_1)
			           +(g_0^{\otimes N}-g_1^{\otimes N})\otimes(g_0\mp g_1)
			           \Bigr].
		\notag
	\end{align}
	Both local combinations are positive semidefinite: each is block
	diagonal in the flag bit, with $C_0+C_1=\frac74P_+$ in one flag sector
	and $C_0-C_1=\frac14P_-$ in the other, the two sectors being exchanged
	between the sum and the difference. A Kronecker product of positive
	semidefinite matrices is positive semidefinite, and so are sums,
	nonnegative multiples, and pullbacks along an arbitrary reindexing.
\end{proof}

% Project issue #7776, not a printed CPSV16 result.
% Provenance: docs/audits/2026-09-05_twisted_dimer_unitary_factorization.md.
\begin{definition}[Controlled Bell-sign flip]%
	\label{def:mpdo_twisted_dimer_controlled_flip}
	\lean{MPOTensor.TwistedDimer.bellProjector, MPOTensor.TwistedDimer.bondState}
	\lean{MPOTensor.TwistedDimer.sigma, MPOTensor.TwistedDimer.sigma'}
	\lean{MPOTensor.TwistedDimer.flagMatrix, MPOTensor.TwistedDimer.flagFlip}
	\lean{MPOTensor.TwistedDimer.localV}
	\leanok
	On two bond qubits let
	\begin{align}
		\ket{\Phi_\pm} & =\frac{\ket{00}\pm\ket{11}}{\sqrt2},
		\notag\\
		\Pi_\pm & =\ketbra{\Phi_\pm}{\Phi_\pm}.
		\notag
	\end{align}
	Define $\sigma=\frac78\Pi_++\frac18\Pi_-$ and
	$\sigma'=\frac78\Pi_+-\frac18\Pi_-$.
	On a flag qubit let $X\ket{f}=\ket{f+1}$, with addition modulo two,
	and $Z\ket{f}=(-1)^f\ket{f}$. The controlled Bell-sign flip is
	\begin{align}
		V & =(\Id-\Pi_-)\otimes\Id+\Pi_-\otimes X.
		\notag
	\end{align}
\end{definition}

\begin{lemma}[Bond-state matrix entries]%
	\label{lem:mpdo_twisted_dimer_bond_state_entries}
	\lean{MPOTensor.TwistedDimer.bondState_apply}
	\leanok
	\uses{def:mpdo_twisted_dimer_tensor, def:mpdo_twisted_dimer_controlled_flip}
	For $k\in\{0,1\}$, put $\sigma_k=\frac78\Pi_++(-1)^k\frac18\Pi_-$,
	where $\Pi_\pm$ are the Bell projections of
	Definition~\ref{def:mpdo_twisted_dimer_controlled_flip}.
	Let $C_k$ be the matrix of Definition~\ref{def:mpdo_twisted_dimer_tensor}.
	For $r,l,r',l'\in\{0,1\}$, in right-then-left bond order,
	\begin{align}
		\bra{r,l}\sigma_k\ket{r',l'} & =
		\begin{cases}
			C_k[l,l'] & \text{if } r=l \text{ and } r'=l', \\
			0         & \text{otherwise}.
		\end{cases}
		\notag
	\end{align}
\end{lemma}
\begin{proof}\leanok
	\uses{def:mpdo_twisted_dimer_tensor, def:mpdo_twisted_dimer_controlled_flip}
	Expand the two Bell projections in the computational basis.
\end{proof}

\begin{lemma}[Local unitary identities]%
	\label{lem:mpdo_twisted_dimer_controlled_flip}
	\lean{MPOTensor.TwistedDimer.localV_conjTranspose}
	\lean{MPOTensor.TwistedDimer.localV_mul_self}
	\lean{MPOTensor.TwistedDimer.localV_conjugate}
	\leanok
	\uses{def:mpdo_twisted_dimer_controlled_flip}
	The operator $V$ of Definition~\ref{def:mpdo_twisted_dimer_controlled_flip}
	is self-adjoint and satisfies $V^2=\Id$ on the full bond--flag space.
	It obeys
	\begin{align}
		V(\sigma\otimes\Id)V^\dagger & =\sigma\otimes\Id,
		\notag\\
		V(\sigma\otimes Z)V^\dagger & =\sigma'\otimes Z.
		\notag
	\end{align}
\end{lemma}
\begin{proof}\leanok
	\uses{def:mpdo_twisted_dimer_controlled_flip}
	Evaluate the $8\times8$ matrices entry by entry in the computational
	basis. This gives $V^\dagger=V$, $V^2=\Id$,
	$V(\sigma\otimes\Id)=(\sigma\otimes\Id)V$, and
	$V(\sigma\otimes Z)=(\sigma'\otimes Z)V$.
	Multiplying each intertwining identity on the right by $V^\dagger$
	gives the stated conjugation identities.
\end{proof}

\begin{definition}[Bond--flag regrouping]%
	\label{def:mpdo_twisted_dimer_bond_flag_regrouping}
	\lean{MPOTensor.TwistedDimer.incomingCellEquiv, MPOTensor.TwistedDimer.bondFlagEquiv}
	\lean{MPOTensor.TwistedDimer.powN, MPOTensor.TwistedDimer.evenFlagState}
	\lean{MPOTensor.TwistedDimer.decoratedState, MPOTensor.TwistedDimer.chainUnitary}
	\lean{MPOTensor.TwistedDimer.chainUnitary_apply}
	\leanok
	\uses{def:mpdo_twisted_dimer_tensor, def:mpdo_twisted_dimer_controlled_flip}
	For $N\ge0$, write each physical configuration as
	$s=((l_m,r_m,f_m))_m$, with cyclic indices when $N>0$.
	At $N=0$ the configurations are empty. The incoming cells are
	$e_N(s)_m=((r_{m-1},l_m),f_m)$.
	For cells $a_m=(b_m,g_m)$, with $b_m\in\{0,1\}^2$, the inverse
	reconstructs site $m$ as $((b_m)_2,(b_{m+1})_1,g_m)$.
	Let $E_N\ket{s}=\ket{e_N(s)}$, and let $S_N$ separate a cell
	configuration into its bond string and flag string. Define
	\begin{align}
		\omega_N & =2^{-N}(\Id^{\otimes N}+Z^{\otimes N}),
		\notag\\
		D_N & =E_N^\dagger S_N^\dagger
		(\sigma^{\otimes N}\otimes\omega_N)S_NE_N,
		\notag\\
		U_N & =E_N^\dagger V^{\otimes N}E_N.
		\notag
	\end{align}
	For $N>0$, $\omega_N$ is the uniform state on even-parity flag strings,
	and $D_N$ is the product of independent bond states and that flag
	state, expressed in the physical site coordinates.
	At length zero, the empty-product convention gives
	$\omega_0=D_0=2\Id$ and $U_0=\Id$.
\end{definition}

\begin{theorem}[Unitary bond--flag factorization]%
	\label{thm:mpdo_twisted_dimer_unitary_factorization}
	\lean{MPOTensor.TwistedDimer.chainUnitary_conjTranspose}
	\lean{MPOTensor.TwistedDimer.chainUnitary_mul_conjTranspose}
	\lean{MPOTensor.TwistedDimer.chainUnitary_conjTranspose_mul}
	\lean{MPOTensor.TwistedDimer.mpo_eq_unitary_factorization}
	\leanok
	\discussion{7776}
	\uses{def:mpdo_twisted_dimer_bond_flag_regrouping}
	For every positive length $N$, the operator $U_N$ of
	Definition~\ref{def:mpdo_twisted_dimer_bond_flag_regrouping} is
	self-adjoint and obeys $U_NU_N^\dagger=U_N^\dagger U_N=\Id$.
	The closed operator of the twisted quantum dimer satisfies
	$\rho^{(N)}(M)=U_ND_NU_N^\dagger$.
\end{theorem}
\begin{proof}\leanok
	\uses{lem:mpdo_twisted_dimer_controlled_flip, lem:mpdo_twisted_dimer_bond_state_entries, thm:mpdo_twisted_dimer_ismpdo}
	Take tensor products of the identities in
	Lemma~\ref{lem:mpdo_twisted_dimer_controlled_flip}. After separating
	the bond and flag strings, we obtain
	\begin{align}
		S_NE_NU_ND_NU_N^\dagger E_N^\dagger S_N^\dagger
		 & =2^{-N}\left(\sigma^{\otimes N}\otimes\Id^{\otimes N}
		          +(\sigma')^{\otimes N}\otimes Z^{\otimes N}\right).
		\notag
	\end{align}
	By Lemma~\ref{lem:mpdo_twisted_dimer_bond_state_entries}, both bond
	operators vanish outside the span of $\ket{00}$ and
	$\ket{11}$; on this span their matrices are $C_0$ and $C_1$,
	respectively. In physical coordinates the entries therefore equal
	the cyclic bond-matching formula of
	Theorem~\ref{thm:mpdo_twisted_dimer_ismpdo}.
\end{proof}

This factorization concerns the closed operators of the present example;
it is not a result stated in \cite{Cirac2016MPDO_arXiv}. Its gates act on
an incoming bond and its receiving flag, and hence cross neighboring
physical sites. It asserts neither an on-site tensor factorization nor a
virtual gauge equivalence. This theorem alone gives neither a
circuit-depth bound nor a classification of the factors as RFP tensors.
Positive length is necessary: the empty closed operator equals $8$,
while the empty product expression equals $2$.

% Project issue #7782; CPSV16 Example 4.12, source lines 932--939,
% with the density normalization already specified in the cited theorem.
\begin{theorem}[Identification of the flag factor]%
	\label{thm:mpdo_twisted_dimer_flag_factor}
	\lean{MPOTensor.TwistedDimer.flagMatrix_zero, MPOTensor.TwistedDimer.flagMatrix_one}
	\lean{MPOTensor.TwistedDimer.evenFlagState_eq_mpo_Mhat}
	\lean{MPOTensor.TwistedDimer.evenFlagState_has_rfpRepresentation}
	\lean{MPOTensor.TwistedDimer.evenFlagState_posSemidef}
	\lean{MPOTensor.TwistedDimer.trace_evenFlagState}
	\leanok
	\uses{def:mpdo_twisted_dimer_bond_flag_regrouping, thm:cpsv_example412_normalized_rfp_maps}
	Let $\widehat M$ be the density-normalized tensor of
	Theorem~\ref{thm:cpsv_example412_normalized_rfp_maps}. For every
	$N\geq0$, the flag operator satisfies
	$\omega_N=\rho^{(N)}(\widehat M)$.
	In particular, the flag family has an MPDO representation satisfying
	both RFP channel equations. For $N>0$, $\omega_N$ is positive
	semidefinite and has trace one. At $N=0$, both
	$\omega_0$ and $\rho^{(0)}(\widehat M)$ are the scalar $2$.
\end{theorem}
\begin{proof}\leanok
	\uses{thm:cpsv_example412_normalized_rfp_maps, thm:cpsv_example412_literal_formula}
	Substitute the identity and Pauli $Z$ matrices in the definition of
	$\omega_N$. Its expression $2^{-N}(\Id^{\otimes N}+Z^{\otimes N})$
	is the closed operator of $\widehat M$ by
	Theorem~\ref{thm:cpsv_example412_literal_formula} and the factor
	$2^{-N}$; the empty case is evaluated directly. Positivity,
	normalization, and the two channel equations
	follow from Theorem~\ref{thm:cpsv_example412_normalized_rfp_maps}.
\end{proof}

\begin{theorem}[Normalization of the bond--flag factorization]%
	\label{thm:mpdo_twisted_dimer_factor_normalization}
	\lean{MPOTensor.TwistedDimer.bellProjector_posSemidef}
	\lean{MPOTensor.TwistedDimer.sigma_posSemidef, MPOTensor.TwistedDimer.trace_sigma}
	\lean{MPOTensor.TwistedDimer.trace_powN_sigma}
	\lean{MPOTensor.TwistedDimer.decoratedState_posSemidef}
	\lean{MPOTensor.TwistedDimer.trace_decoratedState}
	\lean{MPOTensor.TwistedDimer.trace_mpo_T}
	\leanok
	\uses{def:mpdo_twisted_dimer_controlled_flip, def:mpdo_twisted_dimer_bond_flag_regrouping, def:mpdo_twisted_dimer_tensor}
	The bond operator $\sigma$ is positive semidefinite and has trace one.
	Every tensor power $\sigma^{\otimes N}$ has trace one, including
	$N=0$. For every $N>0$, $D_N$ is positive semidefinite and
	$\tr D_N=\tr\rho^{(N)}(M)=1$.
\end{theorem}
\begin{proof}\leanok
	\uses{lem:mpdo_twisted_dimer_bond_state_entries, thm:mpdo_twisted_dimer_flag_factor, thm:mpdo_twisted_dimer_unitary_factorization}
	The Bell projectors are positive rank-one operators. By
	Lemma~\ref{lem:mpdo_twisted_dimer_bond_state_entries},
	$\tr\sigma=C_0[0,0]+C_0[1,1]=1$.
	Tensor products preserve positivity, and their traces multiply.
	Unitary changes of coordinates preserve both positivity and trace.
	Apply these facts to $D_N$ and then use
	Theorem~\ref{thm:mpdo_twisted_dimer_unitary_factorization}.
\end{proof}

The flag identification uses the density-normalized tensor $\widehat M$,
not its unnormalized counterpart. It does not classify the structure
coefficients of a normalized vertical BNT presentation. In particular,
constant multiplication coefficients for the unnormalized representatives
$\Id,Z$ should not be confused with those for the spectrally normalized
representatives $\Id/\sqrt2,Z/\sqrt2$.

% Project issue #7783; round44 construction note, lines 162--193,
% prop:p6-r44-dimer. Only its bond-state and channel construction is used.
\begin{definition}[The mixed-Bell bond dimer]%
	\label{def:mpdo_twisted_dimer_bond_tensor}
	\lean{MPOTensor.TwistedDimer.bondSiteEquiv, MPOTensor.TwistedDimer.sigmaDimer}
	\lean{MPOTensor.TwistedDimer.incomingBondEquiv}
	\leanok
	\uses{def:mpdo_twisted_dimer_controlled_flip, def:mpdo_twisted_dimer_tensor}
	Let
	\begin{align}
		C & =\begin{pmatrix}1/2&3/8\\3/8&1/2\end{pmatrix}.
		\notag
	\end{align}
	A physical site consists of qubits $(l,r)\in\{0,1\}^2$.
	With virtual indices in $\{0,1\}^2$, define the tensor $Q$ by
	\begin{align}
		Q^{(l,r),(l',r')} & =C_{ll'}E_{(l,l'),(r,r')},
		\notag
	\end{align}
	where $E_{a,b}$ is the matrix unit with row $a$ and column $b$.
	Both physical and virtual dimensions are four. For $N\ge0$, let
	$B_N$ be the permutation of basis configurations given by
	\begin{align}
		B_N\ket{((l_m,r_m))_m} & =\ket{((r_{m-1},l_m))_m},
		\notag
	\end{align}
	with cyclic indices when $N>0$ and the unique empty configuration
	when $N=0$. Its inverse sends a bond configuration $(b_m)_m$
	to the site configuration $(((b_m)_2,(b_{m+1})_1))_m$.
\end{definition}

\begin{theorem}[Closed operators of the mixed-Bell dimer]%
	\label{thm:mpdo_twisted_dimer_bond_product}
	\lean{MPOTensor.TwistedDimer.mpo_sigmaDimer_eq_bondProduct}
	\lean{MPOTensor.TwistedDimer.sigmaDimer_isMPDO}
	\lean{MPOTensor.TwistedDimer.trace_mpo_sigmaDimer}
	\leanok
	\uses{def:mpdo_twisted_dimer_bond_tensor, def:mpdo}
	For every $N>0$, the tensor $Q$ of
	Definition~\ref{def:mpdo_twisted_dimer_bond_tensor} satisfies
	\begin{align}
		\rho^{(N)}(Q) & =B_N^\dagger\sigma^{\otimes N}B_N.
		\notag
	\end{align}
	This operator is positive semidefinite and has trace one. In
	particular, $Q$ generates MPDOs. The identity excludes $N=0$, where the
	empty closed operator is $4$ and the empty bond product is $1$.
\end{theorem}
\begin{proof}\leanok
	\uses{def:mpdo_twisted_dimer_bond_tensor, lem:matrix_cyclic_path_sum, lem:mpdo_twisted_dimer_bond_state_entries, thm:mpdo_twisted_dimer_factor_normalization}
	Fix independent configurations $s_m=(l_m,r_m)$ and
	$t_m=(l'_m,r'_m)$, with $m=0,\ldots,N-1$ and cyclic successors.
	The cyclic trace expansion of Lemma~\ref{lem:matrix_cyclic_path_sum}
	has only one possibly nonzero virtual path, namely $(l_m,l'_m)_m$.
	Hence
	\begin{align}
		\tr\!\left(\prod_{m=0}^{N-1}Q^{(l_m,r_m),(l'_m,r'_m)}\right)
		 & =\begin{cases}
			    \prod_{m=0}^{N-1}C_{l_m l'_m}
			      & \text{if } r_m=l_{m+1} \text{ and } r'_m=l'_{m+1}
			    \text{ for every }m,                                  \\
			    0 & \text{otherwise}.
		    \end{cases}
		\label{eq:mpdo_dimer_closed_entries}
	\end{align}
	By Lemma~\ref{lem:mpdo_twisted_dimer_bond_state_entries}, the
	$(s,t)$ entry of $B_N^\dagger\sigma^{\otimes N}B_N$ is
	$\prod_m\sigma_{(r_{m-1},l_m),(r'_{m-1},l'_m)}$.
	Its matching conditions and value agree with
	(\ref{eq:mpdo_dimer_closed_entries}) after cyclic relabeling.
	Positivity and trace one follow from
	Theorem~\ref{thm:mpdo_twisted_dimer_factor_normalization}.
\end{proof}

\begin{definition}[Bond insertion and removal]%
	\label{def:mpdo_twisted_dimer_bond_channels}
	\lean{MPOTensor.TwistedDimer.bondInsertionEquiv}
	\lean{MPOTensor.TwistedDimer.sigmaDimerRefine, MPOTensor.TwistedDimer.sigmaDimerCoarse}
	\leanok
	\uses{def:mpdo_twisted_dimer_controlled_flip, def:mpdo_twisted_dimer_bond_tensor}
	Let $J$ be the basis permutation from an exterior site and an
	inserted bond to two sites:
	\begin{align}
		J\ket{(l,r),(u,v)} & =\ket{(l,u),(v,r)}.
		\notag
	\end{align}
	Define the refinement and coarse-graining maps by
	\begin{align}
		\mathcal T_Q(Y) & =J(Y\otimes\sigma)J^\dagger,
		\notag\\
		\mathcal S_Q(W) & =\tr_{\mathrm{bond}}(J^\dagger WJ).
		\notag
	\end{align}
	The partial trace is over the inserted bond, leaving the two
	exterior qubits as a single physical site.
\end{definition}

\begin{theorem}[Channel equations for the mixed-Bell dimer]%
	\label{thm:mpdo_twisted_dimer_bond_rfp}
	\lean{MPOTensor.TwistedDimer.sigmaDimer_physClose1, MPOTensor.TwistedDimer.sigmaDimer_mul}
	\lean{MPOTensor.TwistedDimer.sigmaDimerRefine_isKrausCPTP}
	\lean{MPOTensor.TwistedDimer.sigmaDimerCoarse_isKrausCPTP}
	\lean{MPOTensor.TwistedDimer.sigmaDimerCoarse_refine}
	\lean{MPOTensor.TwistedDimer.sigmaDimerRefine_physClose1}
	\lean{MPOTensor.TwistedDimer.sigmaDimerCoarse_physClose2}
	\lean{MPOTensor.TwistedDimer.sigmaDimer_isRFPViaTS}
	\leanok
	\uses{def:mpdo_twisted_dimer_bond_channels, def:is_rfp_via_ts}
	The maps $\mathcal T_Q$ and $\mathcal S_Q$ of
	Definition~\ref{def:mpdo_twisted_dimer_bond_channels} are tpCPMs.
	For every one-site operator $Y$,
	$\mathcal S_Q(\mathcal T_Q(Y))=Y$.
	For every virtual boundary matrix $X\in\MN{4}$, the
	one-site and two-site physical closures satisfy
	\begin{align}
		\mathcal T_Q(Q_1(X)) & =Q_2(X),
		\notag\\
		\mathcal S_Q(Q_2(X)) & =Q_1(X).
		\notag
	\end{align}
	Thus $Q$ is an RFP in the sense of
	Definition~\ref{def:is_rfp_via_ts}.
\end{theorem}
\begin{proof}\leanok
	\uses{thm:mpdo_twisted_dimer_factor_normalization, def:mpdo_twisted_dimer_bond_tensor, lem:mpdo_twisted_dimer_bond_state_entries}
	Preparation of the positive trace-one state $\sigma$, permutation
	of basis coordinates, and partial trace are tpCPMs.
	Tracing out the prepared state proves
	$\mathcal S_Q(\mathcal T_Q(Y))=Y$.
	For the refinement equation, fix independent ket and bra indices
	$l,r,l',r',u,v,u',v'\in\{0,1\}$. Multiplication of matrix units gives
	\begin{align}
		Q^{(l,r),(l',r')}Q^{(u,v),(u',v')}
		 & =\begin{cases}
			    C_{ll'}C_{uu'}E_{(l,l'),(v,v')}
			      & \text{if } r=u \text{ and } r'=u', \\
			    0 & \text{otherwise}.
		    \end{cases}
		\notag
	\end{align}
	For every boundary matrix $X\in\MN{4}$, the one-site
	closure has entries $[Q_1(X)]_{(l,r),(l',r')}=C_{ll'}X_{(r,r'),(l,l')}$.
	Put $a=((l,r),(u,v))$ and $b=((l',r'),(u',v'))$.
	The definition of the insertion map and
	Lemma~\ref{lem:mpdo_twisted_dimer_bond_state_entries} give
	\begin{align}
		[\mathcal T_Q(Q_1(X))]_{a,b}
		 & =[Q_1(X)]_{(l,v),(l',v')}\sigma_{(r,u),(r',u')}
		\notag                                             \\
		 & =\begin{cases}
			    C_{ll'}C_{uu'}X_{(v,v'),(l,l')}
			      & \text{if } r=u \text{ and } r'=u', \\
			    0 & \text{otherwise}
		    \end{cases}
		\notag                                             \\
		 & =[Q_2(X)]_{a,b}.
		\notag
	\end{align}
	Thus the refinement equation holds for arbitrary boundary matrices.
	Applying $\mathcal S_Q\mathcal T_Q=\id$ gives the coarse-graining equation.
\end{proof}

These channel equations prove that the mixed-Bell bond dimer is an RFP;
this does not follow merely from its closed-operator factorization.
The nonzero bond-state eigenvalues are $7/8$ and $1/8$.
No simplicity, CF normalization, or structure-coefficient formula for
a vertical BNT presentation of $Q$ is asserted here.

\begin{definition}[On-site separation of bond qubits and flags]%
	\label{def:mpdo_twisted_dimer_onsite_split}
	\lean{MPOTensor.TwistedDimer.onsiteBondFlagEquiv}
	\leanok
	\uses{def:mpdo_twisted_dimer_tensor, def:mpdo_twisted_dimer_bond_tensor}
	For $N$ sites, let $H_N$ be the basis permutation separating each
	site's two bond qubits from its flag:
	\begin{align}
		H_N\ket{((l_m,r_m,f_m))_m}
		 & =\ket{((l_m,r_m))_m}\otimes\ket{(f_m)_m}.
		\notag
	\end{align}
	Unlike the incoming-bond regrouping, this separation does not move
	any qubit between physical sites.
\end{definition}

\begin{theorem}[Factorization into two RFP families]%
	\label{thm:mpdo_twisted_dimer_rfp_factor_attachment}
	\lean{MPOTensor.TwistedDimer.onsiteBondFlagEquiv_incoming}
	\lean{MPOTensor.TwistedDimer.decoratedState_eq_mpo_factors}
	\lean{MPOTensor.TwistedDimer.mpo_eq_unitary_mpo_factors}
	\leanok
	\uses{def:mpdo_twisted_dimer_onsite_split, def:mpdo_twisted_dimer_bond_flag_regrouping, thm:mpdo_twisted_dimer_flag_factor}
	Let $Q$ be the mixed-Bell bond tensor and let $\widehat M$ be the
	density-normalized tensor identified in
	Theorem~\ref{thm:mpdo_twisted_dimer_flag_factor}. For every $N>0$,
	\begin{align}
		D_N & =H_N^\dagger
		(\rho^{(N)}(Q)\otimes\rho^{(N)}(\widehat M))H_N,
		\notag\\
		\rho^{(N)}(M) & =U_NH_N^\dagger
		(\rho^{(N)}(Q)\otimes\rho^{(N)}(\widehat M))
		H_NU_N^\dagger.
		\notag
	\end{align}
	Both factor tensors generate MPDOs and are RFPs, by
	Theorems~\ref{thm:mpdo_twisted_dimer_bond_product},
	\ref{thm:mpdo_twisted_dimer_bond_rfp}, and
	\ref{thm:mpdo_twisted_dimer_flag_factor}.
\end{theorem}
\begin{proof}\leanok
	\uses{thm:mpdo_twisted_dimer_bond_product, thm:mpdo_twisted_dimer_flag_factor, thm:mpdo_twisted_dimer_unitary_factorization}
	The coordinate permutations satisfy
	$(B_N\otimes\Id)H_N=S_NE_N$: both read the incoming bonds
	$(r_{m-1},l_m)$ and the unchanged flags $f_m$.
	Substitute the bond-product formula and the flag identification
	into the definition of $D_N$. The formula for $\rho^{(N)}(M)$ then
	follows from Theorem~\ref{thm:mpdo_twisted_dimer_unitary_factorization}.
	The RFP channel equations for the two factors were proved independently
	in the cited theorems.
\end{proof}

The separation $H_N$ is on-site, but the conjugating operator $U_N$ acts
across neighboring sites. These factorizations therefore do not
assert an on-site or virtual-gauge equivalence of the original tensors.
At $N=0$, $D_0=2$, whereas the product of the two empty MPOs is
$4\cdot2=8$, so the factorization requires positive length.

\begin{theorem}[The twisted quantum dimer is a renormalization fixed point]%
	\label{thm:mpdo_twisted_dimer_is_rfp_via_ts}
	\lean{MPOTensor.TwistedDimer.refineMap}
	\lean{MPOTensor.TwistedDimer.refineMap_isKrausCPTP}
	\lean{MPOTensor.TwistedDimer.refineMap_physClose1}
	\lean{MPOTensor.TwistedDimer.coarseMap}
	\lean{MPOTensor.TwistedDimer.coarseMap_isKrausCPTP}
	\lean{MPOTensor.TwistedDimer.coarseMap_physClose2}
	\lean{MPOTensor.TwistedDimer.isRFPViaTS_T}
	\leanok
	\uses{def:is_rfp_via_ts, def:mpdo_twisted_dimer_tensor}
	The twisted quantum dimer of
	Definition~\ref{def:mpdo_twisted_dimer_tensor} is an RFP in the sense of
	Definition~\ref{def:is_rfp_via_ts}. The refinement tpCPM $\mathcal T$
	has eight Kraus operators, labelled by the two prepared site flags
	$f_1,f_2$ and a bond label $\varepsilon$. It prepares those two flags
	on the output sites together with the bond state
	$(\ket{00}+(-1)^\varepsilon\ket{11})/\sqrt2$ on their shared bond,
	with weight $x$ for $\varepsilon=0$ and $y$ for $\varepsilon=1$,
	and is supported on the incoming flag $f=f_1+f_2+\varepsilon$.
	Flag addition is modulo two. The coarse-graining tpCPM $\mathcal S$
	has sixteen Kraus operators, labelled by a member of the orthonormal
	bond basis $(\ket{00}\pm\ket{11})/\sqrt2$, $\ket{01}$, $\ket{10}$
	together with the two site flags. It measures the shared bond pair
	in that basis, reads the two flags, and returns the outer site qubits
	with flag $f_1+f_2+s$, where $s$ is one for the antisymmetric bond
	outcome and zero otherwise. These maps satisfy
	$\mathcal T[M_1(X)]=M_2(X)$ and $\mathcal S[M_2(X)]=M_1(X)$
	for every bond matrix $X$.
	This example is motivated by the question of
	\cite[lines~995--1010]{Cirac2016MPDO_arXiv}; it is not stated there.
\end{theorem}
\begin{proof}\leanok
	\uses{def:mpdo_twisted_dimer_tensor, thm:mpdo_twisted_dimer_closures}
	Fix a one-site basis index $(l,r,f)$. The two-site indices in the
	support of a refinement Kraus operator carry the outer bits $l$ and
	$r$, the two prepared flags, and a free shared bond bit. Each of the
	four labels with $f_1+f_2+\varepsilon=f$ therefore contributes twice
	the squared amplitude $w_\varepsilon/4$, where $w_0=x$ and $w_1=y$.
	Two of the four labels carry each value of $\varepsilon$, and the
	total is $x+y=1$. Supports belonging to distinct one-site basis
	indices are disjoint, so the off-diagonal entries vanish and the
	Kraus operators resolve the identity.
	For the intertwining relation, the flag signs multiply as
	$(\tau_k)_{f_1f_1}(\tau_k)_{f_2f_2}=(\tau_k)_{gg}$ with $g=f_1+f_2$.
	For shared bond bits $b,b'$, the two prepared bond states give
	\begin{align}
		\frac{x}{4}+(-1)^{b+b'}\frac{y}{4} & =\frac12\,C_0[b,b'],
		\notag\\
		\frac{x}{4}-(-1)^{b+b'}\frac{y}{4} & =\frac12\,C_1[b,b'].
		\notag
	\end{align}
	These are exactly the coefficients in the two-site closure of
	Theorem~\ref{thm:mpdo_twisted_dimer_closures}.

	For the coarse-graining map, write $B_v(b_1,b_2)$ for the amplitude of
	the bond-basis vector $v$ at $(b_1,b_2)$. Its orthonormality gives
	\begin{align}
		\sum_v B_v(b_1,b_2)B_v(c_1,c_2)
		 & =\delta_{b_1c_1}\delta_{b_2c_2}.
		\notag
	\end{align}
	The sum over the two site flags is
	\begin{align}
		\sum_{f_1,f_2}
		\delta_{f_1F_1}\delta_{f_2F_2}
		\delta_{f_1F'_1}\delta_{f_2F'_2}
		 & =\delta_{F_1F'_1}\delta_{F_2F'_2}.
		\notag
	\end{align}
	Orthonormality of the bond basis and the flag Kronecker deltas
	identify all six bits of the two-site indices and give the resolution
	of the identity. On the support of the two-site closure, the two bond
	qubits of each argument agree, so the outcomes $\ket{01}$ and
	$\ket{10}$ contribute nothing. The symmetric and antisymmetric Bell
	outcomes contribute $\frac12(\pm1)^{b+b'}$.
	Summing over the two flag pairs compatible with the output flag
	leaves, for each block label $k$, the identity
	\begin{align}
		\frac12\Bigl(\sum_{b,b'}C_k[b,b']
		       +(\tau_k)_{11}\sum_{b,b'}(-1)^{b+b'}C_k[b,b']\Bigr) & =1.
		\notag
	\end{align}
	The left side is $\frac12(\frac74+\frac14)=1$ in both cases:
	for $k=0$ the alternating sum is $\frac14$ and the flag sign is $1$,
	and for $k=1$ it is $-\frac14$ and the flag sign is $-1$.
	The remaining factor is the one-site coefficient
	$\frac12\,C_k[l,l']\,(\tau_k)_{ff} \delta_{ff'}$ of
	Theorem~\ref{thm:mpdo_twisted_dimer_closures}.
\end{proof}
